\documentclass{article}
\usepackage[utf8]{inputenc}
\usepackage{amsmath}

\begin{document}

\section*{2025 MIT Science Bowl High School Invitational}
\subsection*{Round 13}

\section*{TOSS-UP}
\subsection*{1) BIOLOGY}
\textbf{Short Answer}
In order to translocate electrons from cytosolic NADH made during glycolysis to the electron transport chain, the import of malate into the mitochondrial matrix is coupled with the export of what intermediate formed from the transamination (read: trans-am-uh-NAY-shun) of oxaloacetate (read: ox-ah-low-ASS-uh-tayt)?

\textbf{ANSWER: ASPARTATE}

\section*{VISUAL BONUS}
\subsection*{1) BIOLOGY}
\textbf{Short Answer} Your next bonus is visual. Shown in the image is the common garden phlox (read: flox).
Answer the following two questions: 
1) To what order of insects do the pollinators of this flower belong to?
2) This plant would be expected to produce pollen with how many pores?

\textbf{ANSWER: 1) LEPIDOPTERA; 2) 3}

\section*{TOSS-UP}
\subsection*{2) MATH}
\textbf{Short Answer} Tim rolls a fair tetrahedral die with the numbers 1, 2, 3, and 4 painted on its sides, then calculates the sum of the numbers on the three sides that are visible to him.
What is the expected value of his sum?

\textbf{ANSWER: 7.5}

\section*{BONUS}
\subsection*{2) MATH}
\textbf{Short Answer} How many ordered pairs of positive integers have greatest common divisor 90 and least common multiple 1800?

\textbf{ANSWER: 4}

\section*{TOSS-UP}
\subsection*{3) ENERGY}
\textbf{Short Answer} Physicists at the Laboratory for Nuclear Science used spectroscopy to study the normal modes of radium monofluoride, a diatomic molecule.
At room temperature, how many active normal modes does radium monofluoride have?

\textbf{ANSWER: 5}

\section*{VISUAL BONUS}
\subsection*{3) ENERGY}
\textbf{Short Answer} Your next bonus is visual. Researchers at the Koch Institute have been studying tumor metabolism.
Displayed is an MRSI image of a squamous cell carcinoma, which integrates PET (read: pet) and MRI (read: MRI).
Answer the following two questions:
1) What effect is responsible for the increased production of lactate in the tumor compared to the normal cells?
2) In the images, what does FDG stand for?

\textbf{ANSWER: 1) WARBURG EFFECT; 2) FLUORODEOXYGLUCOSE (ACCEPT: FLUOROGLUCOSE, FDG).}

\section*{TOSS-UP}
\subsection*{4) PHYSICS}
\textbf{Short Answer} Rank the following three particles in terms of increasing rest mass: 1) Tau neutrino; 2) Muon; 3) W boson.

\textbf{ANSWER: 1, 2, 3}

\section*{BONUS}
\subsection*{4) PHYSICS}
\textbf{Multiple Choice} A two-dimensional incompressible fluid is originally at uniform height, with a constant gravitational field pointing in the negative y direction.
A horizontal force is applied to each section of liquid proportional to the height of the section relative to the ground at $y=0$. Which of the following shapes best describes the equilibrium shape of the liquid?
W) Parabolic curve
X) Exponential curve
Y) Catenary curve
Z) Hyperbolic curve

\textbf{ANSWER: X) EXPONENTIAL CURVE}

\section*{TOSS-UP}
\subsection*{5) EARTH AND SPACE}
\textbf{Short Answer} David the meterologist is watching radar imagery of a dry desert region and sees a strong precipitation signal overhead, yet the ground remains completely dry.
The sky appears streaked beneath the clouds but with no impact on local humidity measurements.
What kind of precipitation phenomena is he observing?

\textbf{ANSWER: VIRGA (ACCEPT: DRY STORM)}

\section*{BONUS}
\subsection*{5) EARTH AND SPACE}
\textbf{Short Answer} In contrast to the r process and s process, which process is hypothesized to occur in accreting binary systems with neutron stars that occurs on the proton-rich side of stability instead of the neutron side of stability?

\textbf{ANSWER: RP-PROCESS (ACCEPT: RAPID PROTON CAPTURE PROCESS)}

\section*{6) CHEMISTRY}
\subsection*{TOSS-UP}
\textbf{Short Answer} To one significant figure, how many joules of heat are required to raise the temperature of one mole of argon gas at constant pressure by 10 Kelvin?

\textbf{ANSWER: 200}

\section*{BONUS}
\subsection*{6) CHEMISTRY}
\textbf{Multiple Choice} Which of the following are the normalized coefficients of the nitrogen and oxygen 2p orbitals, respectively, in the linear combination describing the highest occupied molecular orbital of nitric oxide?
W) 0.91 and 0.43
X) 0.91 and -0.43 (read: negative 0.43)
Y) 0.43 and 0.91
Z) 0.43 and -0.91 (read: negative 0.91)

\textbf{ANSWER: X) 0.91 AND -0.43 (READ: NEGATIVE 0.43)}

\section*{TOSS-UP}
\subsection*{7) PHYSICS}
\textbf{Multiple Choice} Jolie observes that relativistic charged particles emit electromagnetic radiation when traveling around a black hole.
What type of radiation is Jolie observing?
W) Bremsstrahlung (read: bremz-struh-luhng)
X) Cherenkov
Y) Cyclotron
Z) Synchrotron

\textbf{ANSWER: Z) SYNCHROTRON}

\section*{BONUS}
\subsection*{7) PHYSICS}
\textbf{Short Answer} Identify all of the following three processes that exhibit parity symmetry: 1) Beta decay of carbon-14; 2) Nuclear fission of uranium-238; 3) Ejection of an electron from zinc-64 after absorbing a photon.

\textbf{ANSWER: 2 AND 3}

\section*{TOSS-UP}
\subsection*{8) BIOLOGY}
\textbf{Short Answer} Identify all the following three hormones that are dervived from the same precursor as ACTH: 1) Vasopressin; 2) Melatonin; 3) Melanocyte-stimulating hormone.

\textbf{ANSWER: 3 ONLY}

\section*{VISUAL BONUS}
\subsection*{8) BIOLOGY}
\textbf{Short Answer} Your next bonus is visual. The diagram depicts the relative proportion of a recessive deleterious allele in three populations colored red, blue, and green that are subjected to different conditions, leading to different mutation-selection equillibria.
Identify all of the following three statements which are true: 1) All populations eventually reach fixation; 2) The blue population could be exposed to more intercalating agents than the population in red; 3) The green population could be in Hardy-Weinberg equillibrium.

\textbf{ANSWER: 2 AND 3}

\section*{TOSS-UP}
\subsection*{9) CHEMISTRY}
\textbf{Short Answer} Identify all of the following three functional groups that can be reduced by sodium borohydride: 1) Primary amide; 2) Ester; 3) Ketone.

\textbf{ANSWER: 3 ONLY}

\section*{VISUAL BONUS}
\subsection*{9) CHEMISTRY}
\textbf{Short Answer} Your next bonus is visual. Shown in the image is an important reagent in organic chemistry.
Answer the following two questions: 1) What named reaction is this reagent used in?
2) When an aldehyde is treated with this reagent then, specifying stereochemistry, what functional group is produced?

\textbf{ANSWER: 1) HORNER WADSWORTH EMMONS REACTION (ACCEPT: WITTIG REACTION); 2) E-ALKENE (DO NOT ACCEPT: Z-ALKENE, ALKENE) (ACCEPT: 1) SCHLOSSER MODIFICATION; 2) TRANS ALKENE, DO NOT ACCEPT: 2) ALKENE)}

\section*{TOSS-UP}
\subsection*{10) ENERGY}
\textbf{Short Answer} Chemists from the Coley Group used machine learning to predict the outcome of Diels-Alder reactions.
When buta-1,3-diene reacts with ethene through a Diels-Alder reaction, what is the major product?

\textbf{ANSWER: CYCLOHEXENE (DO NOT ACCEPT: CYCLOHEXANE)}

\section*{BONUS}
\subsection*{10) ENERGY}
\textbf{Multiple Choice} Scientists from the Bosak Group used the presence of cerium in rock samples from a Martian crater as a signal for whether the crater had once been a marine environment.
Which of the following best explains why cerium is a useful signal for marine environments on Mars?
W) Cerium is most commonly found in water-bearing meteorites that formed martian craters
X) Cerium forms highly soluble salts that only precipitate in phosphate-rich marine environments
Y) Cerium readily oxidizes in the presence of water to an easily distinguishable redox state
Z) Cerium can be differentiated from iron oxides in martian soil by Mossbauer spectroscopy

\textbf{ANSWER: Y) CERIUM READILY OXIDIZES IN THE PRESENCE OF WATER TO AN EASILY DISTINGUISHABLE REDOX STATE}

\section*{TOSS-UP}
\subsection*{11) MATH}
\textbf{Multiple Choice} If Tim's bank account compounds 1\% interest each month, which of the following is closest to the total percent increase in the account over the course of 100 months?
W) 100\%
X) 135\%
Y) 170\%
Z) 200\%

\textbf{ANSWER: Y) 170\%}

\section*{BONUS}
\subsection*{11) MATH}
\textbf{Short Answer} A non leap calendar year has 525,600 minutes.
What is the sum of all of the unique prime divisors of 525,600?

\textbf{ANSWER: 83}

\section*{TOSS-UP}
\subsection*{12) EARTH AND SPACE}
\textbf{Short Answer} Order the following three supercontinents from oldest to youngest: 1) Rodinia; 2) Ur; 3) Pangea.

\textbf{ANSWER: 2, 1, 3}

\section*{BONUS}
\subsection*{12) EARTH AND SPACE}
\textbf{Short Answer} Identify all of the following three sites that likely exhibit graded bedding: 1) An ancient river delta; 2) An abyssal plain; 3) A shoreline after a period of transgression.

\textbf{ANSWER: ALL}

\section*{13) MATH}
\subsection*{TOSS-UP}
\textbf{Multiple Choice} Which of the following distributions has a mean strctly greater than its median?
W) Uniform
X) Normal
Y) Binomial
Z) Geometric

\textbf{ANSWER: Z) GEOMETRIC}

\section*{BONUS}
\subsection*{13) MATH}
\textbf{Short Answer} Vectors $a$ and $b$ have magnitudes 6 and 8, respectively, and their dot product is 16. What is the magnitude of the vector $a+b$?

\textbf{ANSWER: $2\sqrt{33}$}

\section*{TOSS-UP}
\subsection*{14) EARTH AND SPACE}
\textbf{Short Answer} Soft gamma ray repeaters are hypothesized to result from the occurrence of starquakes on what specific type of pulsars?

\textbf{ANSWER: MAGNETARS}

\section*{VISUAL BONUS}
\subsection*{14) EARTH AND SPACE}
\textbf{Short Answer} Your next bonus is visual.
Answer the following two questions about this constellation: 1) What is the common name of its brightest star?
2) In what month does that star culminate around midnight?

\textbf{ANSWER: 1) REGULUS; 2) FEBRUARY}

\section*{TOSS-UP}
\subsection*{15) ENERGY}
\textbf{Short Answer} Researchers at CSAIL's Julia Lab recently published a study on discretized Brownian motions with random matrices.
The root-mean square displacemnt traveled by a particle experiencing only Brownian motion is proportional to the time elapsed raised to what power?

\textbf{ANSWER: $1/2$}

\section*{BONUS}
\subsection*{15) ENERGY}
\textbf{Multiple Choice} Researchers at CSAIL's Julia Lab used matrix transforms to accelerate computation of the Hamming distance of two inputs.
Which of the following tasks will this development most greatly affect?
W) Finding shortest path on a geodesic
X) Measuring length of qubit arrays
Y) Calculating similarity between genetic sequences
Z) Recovering information from compressed signals

\textbf{ANSWER: Y) CALCULATING SIMILARITY BETWEEN GENETIC SEQUENCES}

\section*{16) PHYSICS}
\subsection*{TOSS-UP}
\textbf{Multiple Choice} Consider a thin, flat, and positively charged conducting disk.
How does its surface charge density change with distance from the disk's center?
W) Increases rapidly, then approaches a constant
X) Increases slowly, then increases rapidly
Y) Decreases rapidly, then approaches a constant
Z) Decreases slowly, then decreases rapidly

\textbf{ANSWER: X) INCREASES SLOWLY, THEN INCREASES RAPIDLY}

\section*{VISUAL BONUS}
\subsection*{16) PHYSICS}
\textbf{Short Answer} Your next bonus is visual. The given diagram depicts the motion of a planet orbiting a rotating star.
Identify all of the following three quantities that are conserved in this system: 1) Total energy of the planet; 2) Angular momentum of the planet; 3) Semi-major axis of the planet's orbit.

\textbf{ANSWER: NONE}

\section*{TOSS-UP}
\subsection*{17) CHEMISTRY}
\textbf{Short Answer} At high temperatures, the force of attraction between two non-polar ideal gas molecules is proportional to their distance of separation raised to what power?

\textbf{ANSWER: -7}

\section*{BONUS}
\subsection*{17) CHEMISTRY}
\textbf{Multiple Choice} The entropy of a system is proportional to the partial derivative with respect to temperature of which of the following thermodynamic state functions at constant pressure?
W) Internal energy
X) Enthalpy
Y) Gibbs free energy
Z) Hemholtz free energy

\textbf{ANSWER: Y) GIBBS FREE ENERGY}

\section*{TOSS-UP}
\subsection*{18) BIOLOGY}
\textbf{Short Answer} A Winogradsky (read: WIN-no-GRAHD-skee) column is an open tube with mud and water containing a diverse microbial ecosystem.
Assuming they can are able to grow in such a column, order the following three microbes by where they would be found from top to bottom.
1) C. difficile; 2) Cyanobacteria; 3) Yeast.

\textbf{ANSWER: 2, 3, 1}

\section*{VISUAL BONUS}
\subsection*{18) BIOLOGY}
\textbf{Short Answer} Your next bonus is visual. Shown in the image is a Wright-Giemsa (read: rite-geem-sa) stain of a sample of blood cells collected from a patient with cancer.
Answer the following two questions: 1) Identify the cells labeled A and B;
2) The cells known as B undergo what process to release cytotoxic compounds from secretory vesicles?

\textbf{ANSWER: 1) 1 IS EOSINOPHIL, 2 IS LYMPHOCYTE (ACCEPT: B CELL, T CELL) 2) DEGRANULATION}

\section*{TOSS-UP}
\subsection*{19) PHYSICS}
\textbf{Short Answer} When alternating current is run through a conducting wire, current density tends to be most densely distributed near the surface of the wire.
What is the name for this phenomenon?

\textbf{ANSWER: SKIN EFFECT}

\section*{BONUS}
\subsection*{19) PHYSICS}
\textbf{Short Answer} Phoebe launches a projectile at an angle of $\pi/4$ (read: pi over four) radians above the horizontal, causing the projectile to land 100 meters away.
To three significant figures, how far would the projectile have landed if the launch angle was reduced by 0.05 radians?

\textbf{ANSWER: 99.5}

\section*{20) CHEMISTRY}
\subsection*{TOSS-UP}
\textbf{Multiple Choice} Which of the following is the molecular geometry of the central atom in xenon dioxydifluoride (read: di-ok-see-dif-luor-ide)?
W) T-shaped
X) Tetrahedral
Y) Square planar
Z) See-saw

\textbf{ANSWER: Z) SEE-SAW}

\section*{BONUS}
\subsection*{20) CHEMISTRY}
\textbf{Multiple Choice} Which of the following mass to charge ratios corresponds to the peak with the greatest relative intensity in the electron ionization mass spectrum of tert-butyl fluoride?
W) 76
X) 75
Y) 61
Z) 57

\textbf{ANSWER: Y) 61}

\section*{TOSS-UP}
\subsection*{21) MATH}
\textbf{Short Answer} The numbers of vertices, faces, and edges are summed for a polyhedron, and the result is 50. How many edges does this polyhedron have?

\textbf{ANSWER: 24}

\section*{BONUS}
\subsection*{21) MATH}
\textbf{Short Answer} Two Science Bowl players form a team.
Each player uniformly, randomly, and independently specializes in 2 of the 6 science bowl subjects.
What is the expected number of unique subjects specialized in by at least one member of their team?

\textbf{ANSWER: 10/3}

\section*{TOSS-UP}
\subsection*{22) BIOLOGY}
\textbf{Short Answer} Jonathan takes a enzyme called X, reversibly inhibits it 3 ways, and then measures the Km of the enzyme after inhibition.
Order the following three types of inhibition by increasing Km of X: 1) Noncompetitive; 2) Uncompetitive; 3) Competitive.

\textbf{ANSWER: 2, 1, 3}

\section*{BONUS}
\subsection*{22) BIOLOGY}
\textbf{Short Answer} In a particular ecosystem at equilibrium, milk snakes are Batesian (read: bates-ian) mimics of Sonoran coral snakes and Texas coral snakes are Mertensian (read: mer-tense-ian) mimics of Sonoran coral snakes.
Order the following three species by increasing relative population change after a non-native predatory falcon species is introduced to this ecosystem: 1) Milk snake, 2) Sonoran coral snake, 3) Texas coral snake

\textbf{ANSWER: 1, 3, 2}

\section*{TOSS-UP}
\subsection*{23) ENERGY}
\textbf{Short Answer} The Winslow Lab is contributing to KamLAND (read: cam-land), an experiment to study double-beta decay in xenon-136.
Unlike beta decay, the double beta-decay of xenon-136 does not produce what elementary particle?

\textbf{ANSWER: NEUTRINO}

\section*{BONUS}
\subsection*{23) ENERGY}
\textbf{Multiple Choice} The Vuletić Group recently published a study probing fermion dynamics on a quantum computer, where fermions were represented as complex numbers.
Which of the following constraints must hold for the complex number representation of a Majorana fermion?
W) The real and imaginary parts are equal in both magnitude and sign
X) The real and imaginary parts are equal in magnitude but opposite in sign
Y) The real part is zero while the imaginary part is non-zero
Z) The real part is non-zero while the imaginary part is zero

\textbf{ANSWER: Z) THE REAL PART IS NON-ZERO WHILE THE IMAGINARY PART IS ZERO}

\section*{24) EARTH AND SPACE}
\subsection*{TOSS-UP}
\textbf{Multiple Choice} The Encke gap of Saturn's E ring is home to which shepherd moon?
W) Mimas
X) Titan
Y) Pan
Z) Enceladus

\textbf{ANSWER: Y) PAN}

\section*{BONUS}
\subsection*{24) EARTH AND SPACE}
\textbf{Short Answer} The cooling of neutron stars and white dwarfs can be attributed to the emission of neutrinos in what process?

\textbf{ANSWER: URCA PROCESS}

\section*{TOSS-UP}
\subsection*{25) CHEMISTRY}
\textbf{Short Answer} Rank the following three compounds by increasing retention factor in column chromatography [cro-muh-TOG-ruh-fee] with a silica stationary phase and 50:50 mixture of n-hexane and ethyl acetate: 1) 1-hexanol; 2) 2-hexanone; 3) n-Hexane.

\textbf{ANSWER: 1, 2, 3}

\section*{BONUS}
\subsection*{25) CHEMISTRY}
\textbf{Multiple Choice} Which of the following dienophiles would react most quickly with 1,3-butadiene in a Diels-Alder reaction?
W) Ethene
X) Propene
Y) Chloroethene
Z) 2-propenal

\textbf{ANSWER: Z) 2-PROPENAL}

\end{document}